\section{Semantic accuracy before the prescribed certificate}
\label{sec:terminal-modal-semantic-stop}

The logarithmic block in \cref{thm:terminal-modal-log-unconditional} is a
certificate statement.  The explicit all-time nonlinear regime gives a much
earlier semantic upper bound for the continued literal iterate.  This section
keeps those two stopping notions separate.

\begin{theorem}[An explicit constant-time-scale semantic stop]
\label{thm:terminal-modal-semantic-stop}
Set
\begin{equation}
 k_{\rm sem}=\left\lceil\frac{15}{4q}\right\rceil.
 \label{eq:terminal-modal-semantic-stop-time}
\end{equation}
For every $m\ge8192$, the continued named full-face recurrence agrees with
its linear candidate through time $k_{\rm sem}$ and satisfies
\begin{equation}
 \norm{p_{k_{\rm sem}}-p^\star}_\infty<\frac q5=\tau.
 \label{eq:terminal-modal-semantic-stop}
\end{equation}
Consequently the state $p_{k_{\rm sem}}$ has degree-normalized PPR error less
than $\rho+\tau=2q/5$ under the standard RPPR bias bridge.
\end{theorem}

\begin{proof}
Let $k=k_{\rm sem}$ and $s=qk$, so $s\ge15/4$.
\Cref{cor:terminal-modal-static-regime-cutoff} identifies the continued
literal state with the full-face linear candidate and makes its safe envelope
strictly unclipped at every full-face time.  In particular, both facts hold
through time $k$.

The exact error formula
\eqref{eq:terminal-modal-position-candidate-kernel}, the lower entry bounds
\eqref{eq:terminal-modal-entry-position-bounds}, and positivity of $K_k,L,J_k$
with row masses $1,1,k$ give
\begin{equation}
 p_k-p^\star
 \ge-q\delta^k\left(\frac{157}{20}+\frac{s}{8}\right)\one
 \ge-qe^{-s}\left(\frac{157}{20}+\frac{s}{8}\right)\one.
 \label{eq:terminal-modal-semantic-error-lower}
\end{equation}

For the other side, write
\[
 M_k=\max(0,\max_j r_j(p_k)),
 \qquad
 \Delta_k=M_k/q^2.
\]
The unclipped safe envelope is $\ell_k=p_k-\Delta_k\one$.  Since the
normalized full-face operator sends $\one$ to $q^2\one$,
$r(\ell_k)=r(p_k)-M_k\one\le0$.  The full-face operator is a nonsingular
$M$-matrix, so its inverse is entrywise nonnegative; comparison with
$r(p^\star)=0$ yields $\ell_k\le p^\star$.  Because the literal and linear
states agree, $\Delta_k=\delta_k^{\rm lin}$ in the notation of
\eqref{eq:terminal-modal-late-envelope-upper}.  Hence
\begin{equation}
 p_k-p^\star\le\Delta_k\one
 \le\frac{21}{80}qe^{-s}(1+8s)\one,
 \label{eq:terminal-modal-semantic-error-upper}
\end{equation}
where the final inequality is
\eqref{eq:terminal-modal-late-envelope-upper} with the static target supplied
by \cref{cor:terminal-modal-static-regime-cutoff}.

Both scalar factors on the right of
\eqref{eq:terminal-modal-semantic-error-lower}--
\eqref{eq:terminal-modal-semantic-error-upper} decrease for $s\ge15/4$.
The positive exponential series gives the exact lower bound
\begin{equation}
 e^{15/4}>
 \sum_{n=0}^{8}\frac{(15/4)^n}{n!}
 =\frac{2459867141}{58720256}.
 \label{eq:terminal-modal-semantic-exp-fifteen-fourths}
\end{equation}
At $s=15/4$, the two required comparisons reduce to
\begin{align}
 \frac{2459867141}{58720256}-\frac{1331}{32}
 &=\frac{17471493}{58720256}>0,\notag\\
 \frac{2459867141}{58720256}-\frac{651}{16}
 &=\frac{70686725}{58720256}>0.
 \label{eq:terminal-modal-semantic-rational-slacks}
\end{align}
The first proves that the magnitude of the lower bound is less than $q/5$;
the second proves
$(21/80)e^{-15/4}(1+8(15/4))<1/5$.  Thus every coordinate lies strictly
between $-q/5$ and $q/5$, proving
\eqref{eq:terminal-modal-semantic-stop}.  Adding the RPPR bias $\rho=q/5$
gives the PPR statement.
\end{proof}

\begin{corollary}[Semantic-return swept-volume ledger; named family]
\label{cor:terminal-modal-semantic-work}
Run the named execution until the earlier of its prescribed one-sided
certificate and $k_{\rm sem}$.  If the certificate fires first, use its
standard semantic guarantee; otherwise return the theorem-timed state.  The
aggregate face volume of the terminal proximal scans is at most
\begin{equation}
 \vol(P_m)k_{\rm sem}
 \le\frac1{8q}\left(\frac{15}{4q}+1\right)
 =\frac{15}{32q^2}+\frac1{8q}.
 \label{eq:terminal-modal-semantic-terminal-work}
\end{equation}
Including the aggregate face volume of the proper-prefix proximal scans gives
\begin{equation}
 \sum_{j=0}^{m-1}(1+2j)+\vol(P_m)k_{\rm sem}
 \le\frac{121}{256q^2}+\frac1{8q}=O(q^{-2}).
 \label{eq:terminal-modal-semantic-total-work}
\end{equation}
Every named fixed-face step uses only a constant number of passes over its
face.  Hence the complete charged work through this return, including output
and final residual validation, is $O(q^{-2})$.  In the note's PPR namespace
$\epsppr=\rho+\tau=2q/5$, this is $O(\epsppr^{-2})$.
\end{corollary}

\begin{proof}
Use $\vol(P_m)=2m=1/(8q)$ and
$k_{\rm sem}\le15/(4q)+1$.  The prefix volumes sum to
$\sum_{j=0}^{m-1}(1+2j)=m^2=1/(256q^2)$.
As in the proof of the exact-CG escape proposition above, the explicit prefix
optimum and rank-one displacement formulas are materialized by generating
$\chi^i,\chi^{-i}$ sequentially.  Thus response/transport state writes, the
gate, and the prefix update are all absorbed in $O(\vol(U_j))$ work rather
than requiring an uncharged restricted solve.
The named constant-pass implementation charges $O(\vol(U))$ per fixed-face
step, while output and final validation add $O(m)=O(q^{-1})$.
\end{proof}

\begin{remark}[What the semantic theorem does and does not change]
The stopping time in \cref{thm:terminal-modal-semantic-stop} is a
family-specific theorem schedule; it is not the literal implementation's
prescribed one-sided certificate.  If that implementation is required to
wait for its certificate, \cref{cor:terminal-modal-charged-work} still gives
the logarithmically larger charged ledger for all sufficiently large $m$.
Conversely, the semantic theorem shows that this delay cannot be interpreted
as an RPPR accuracy lower bound.  It gives neither an eleven-resource class
lower bound nor a result for other graphs, schedules, or algorithms.
\end{remark}
